\section{R}

% raten (to advise / to guess)
\begin{verbentry}{
  style = {dat},
  toc = {raten (to advise / to guess)},
  name = {raten},
  english = {to advise / to guess},
  type = {irregular, + Dativ},
  auxiliary = {haben}
}
 
    
     \vspace{5pt}

    \hrule
    
    \vspace{5pt}

  \begin{tabular}{rl}
     \multicolumn{2}{l}{\textbf{PRÄSENS} }\\
    ich & \textbf{rate}\\
    du & \textbf{rätst}\\
    er/sie/es & \textbf{rät}\\
    wir & \textbf{raten}\\
    ihr & \textbf{ratet}\\
    sie/Sie & \textbf{raten}\\
  \end{tabular}
  \hfill
     \begin{tabular}{rl}
    \multicolumn{2}{l}{\textbf{PRÄTERITUM} }\\
    ich & \textbf{riet}\\
    du & \textbf{rietest}\\
    er/sie/es & \textbf{riet}\\
    wir & \textbf{rieten}\\
    ihr & \textbf{rietet}\\
    sie/Sie & \textbf{rieten}\\
  \end{tabular}
  \hfill
  \begin{tabular}{rl}
     \multicolumn{2}{l}{\textbf{IMPERATIV} }\\
   & \textbf{rat(e) (du)!}\\
   & \textbf{raten wir!}\\
   & \textbf{ratet (ihr)!}\\
   & \textbf{raten Sie!}\\
   & \vspace{10pt}
  \end{tabular}

    \vspace{10pt}

 \begin{tabular}{rl}
     \multicolumn{2}{l}{\textbf{PERFEKT}}\\
   & haben + \textbf{geraten}\\
  \end{tabular}
\hfill
   \begin{tabular}{rl}
   
    \multicolumn{2}{l}{\textbf{FUTURE I} }\\
   & werden + \textbf{raten}\\
 
  \end{tabular}
\hfill
   \begin{tabular}{rl}
      \multicolumn{2}{l}{\textbf{PLUSQUAMPERFEKT}}\\
   & hatten + \textbf{geraten}\\
  \end{tabular}

   \vspace{10pt}

   \begin{tabular}{lll}
      \multicolumn{3}{l}{\textbf{MIT PRÄPOSITIONEN}}\\
& \textbf{raten zu} + Dat & (to advise to do / recommend)\\
& \textbf{raten von} + Dat & (to advise against)\\
\end{tabular}
  \hfill
   \begin{tabular}{lll}
      \multicolumn{3}{l}{\textbf{TRENNBARE VERBEN}}\\
 & \textbf{ab$|$raten} & (to advise against)\\
 & \textbf{zu$|$raten} & (to advise to do)\\
  \end{tabular}


  \vspace{-5pt}
\end{verbentry}

% reagieren (to react)
\begin{verbentry}{
  style = {acc},
  toc = {reagieren (to react)},
  name = {reagieren},
  english = {to react},
  type = {irregular, + Akkusativ},
  auxiliary = {haben}
}
 
    
     \vspace{5pt}

    \hrule
    
    \vspace{5pt}

  \begin{tabular}{rl}
     \multicolumn{2}{l}{\textbf{PRÄSENS} }\\
    ich & \textbf{reagiere}\\
    du & \textbf{reagierst}\\
    er/sie/es & \textbf{reagiert}\\
    wir & \textbf{reagieren}\\
    ihr & \textbf{reagiert}\\
    sie/Sie & \textbf{reagieren}\\
  \end{tabular}
  \hfill
     \begin{tabular}{rl}
    \multicolumn{2}{l}{\textbf{PRÄTERITUM} }\\
    ich & \textbf{reagierte}\\
    du & \textbf{reagiertest}\\
    er/sie/es & \textbf{reagierte}\\
    wir & \textbf{reagierten}\\
    ihr & \textbf{reagiertet}\\
    sie/Sie & \textbf{reagierten}\\
  \end{tabular}
  \hfill
  \begin{tabular}{rl}
     \multicolumn{2}{l}{\textbf{IMPERATIV} }\\
   & \textbf{reagier(e) (du)!}\\
   & \textbf{reagieren wir!}\\
   & \textbf{reagiert (ihr)!}\\
   & \textbf{reagieren Sie!}\\
   & \vspace{10pt}
  \end{tabular}

    \vspace{10pt}

 \begin{tabular}{rl}
     \multicolumn{2}{l}{\textbf{PERFEKT}}\\
  &  haben + \textbf{reagiert}\\
  \end{tabular}
\hfill
   \begin{tabular}{rl}
   
    \multicolumn{2}{l}{\textbf{FUTURE I} }\\
  &  werden + \textbf{reagieren}\\
 
  \end{tabular}
\hfill
   \begin{tabular}{rl}
      \multicolumn{2}{l}{\textbf{PLUSQUAMPERFEKT}}\\
  &  hatten + \textbf{reagiert}\\
  \end{tabular}
  
  
\vspace{10pt}

\begin{tabular}{lll}
\multicolumn{3}{l}{\textbf{MIT PRÄPOSITIONEN}}\\
& \textbf{---} & \\
\end{tabular}
\hfill
\begin{tabular}{lll}
\multicolumn{3}{l}{\textbf{TRENNBARE VERBEN}}\\
& \textbf{---} & \\
\end{tabular}
\vspace{-5pt}
\end{verbentry}

% reden (to talk)
\begin{verbentry}{
  style = {default},
  toc = {reden (to talk)},
  name = {reden},
  english = {to talk},
  type = {regular},
  auxiliary = {haben}
}
 
    
     \vspace{5pt}

    \hrule
    
    \vspace{5pt}

  \begin{tabular}{rl}
     \multicolumn{2}{l}{\textbf{PRÄSENS} }\\
    ich & \textbf{rede}\\
    du & \textbf{redest}\\
    er/sie/es & \textbf{redet}\\
    wir & \textbf{reden}\\
    ihr & \textbf{redet}\\
    sie/Sie & \textbf{reden}\\
  \end{tabular}
  \hfill
     \begin{tabular}{rl}
    \multicolumn{2}{l}{\textbf{PRÄTERITUM} }\\
    ich & \textbf{redete}\\
    du & \textbf{redetest}\\
    er/sie/es & \textbf{redete}\\
    wir & \textbf{redeten}\\
    ihr & \textbf{redetet}\\
    sie/Sie & \textbf{redeten}\\
  \end{tabular}
  \hfill
  \begin{tabular}{rl}
     \multicolumn{2}{l}{\textbf{IMPERATIV} }\\
   & \textbf{red(e) (du)!}\\
   & \textbf{reden wir!}\\
   & \textbf{redet (ihr)!}\\
   & \textbf{reden Sie!}\\
   & \vspace{10pt}
  \end{tabular}

    \vspace{10pt}

 \begin{tabular}{rl}
     \multicolumn{2}{l}{\textbf{PERFEKT}}\\
  &  haben + \textbf{geredet}\\
  \end{tabular}
\hfill
   \begin{tabular}{rl}
   
    \multicolumn{2}{l}{\textbf{FUTURE I} }\\
  &  werden + \textbf{reden}\\
 
  \end{tabular}
\hfill
   \begin{tabular}{rl}
      \multicolumn{2}{l}{\textbf{PLUSQUAMPERFEKT}}\\
  &  hatten + \textbf{geredet}\\
  \end{tabular}
  
  \vspace{10pt}



\begin{tabular}{lll}
\multicolumn{3}{l}{\textbf{MIT PRÄPOSITIONEN}}\\
& \textbf{---} & \\
\end{tabular}
\hfill
\begin{tabular}{lll}
\multicolumn{3}{l}{\textbf{TRENNBARE VERBEN}}\\
& \textbf{---} & \\
\end{tabular}

  \vspace{-5pt}
\end{verbentry}

% reduzieren (to reduce)
\begin{verbentry}{
  style = {acc},
  toc = {reduzieren (to reduce)},
  name = {reduzieren},
  english = {to reduce},
  type = {regular, + Akkusativ},
  auxiliary = {haben}
}


\vspace{5pt}
\hrule
\vspace{5pt}

\begin{tabular}{rl}
\multicolumn{2}{l}{\textbf{PRÄSENS}}\\
ich & \textbf{reduziere}\\
du & \textbf{reduzierst}\\
er/sie/es & \textbf{reduziert}\\
wir & \textbf{reduzieren}\\
ihr & \textbf{reduziert}\\
sie/Sie & \textbf{reduzieren}\\
\end{tabular}
\hfill
\begin{tabular}{rl}
\multicolumn{2}{l}{\textbf{PRÄTERITUM}}\\
ich & \textbf{reduzierte}\\
du & \textbf{reduziertest}\\
er/sie/es & \textbf{reduzierte}\\
wir & \textbf{reduzierten}\\
ihr & \textbf{reduziertet}\\
sie/Sie & \textbf{reduzierten}\\
\end{tabular}
\hfill
\begin{tabular}{rl}
\multicolumn{2}{l}{\textbf{IMPERATIV}}\\
& \textbf{reduziere (du)!}\\
& \textbf{reduzieren wir!}\\
& \textbf{reduziert (ihr)!}\\
& \textbf{reduzieren Sie!}\\
\end{tabular}

\vspace{10pt}

\begin{tabular}{rl}
\multicolumn{2}{l}{\textbf{PERFEKT}}\\
& haben + \textbf{reduziert}\\
\end{tabular}
\hfill
\begin{tabular}{rl}
\multicolumn{2}{l}{\textbf{FUTURE I}}\\
& werden + \textbf{reduzieren}\\
\end{tabular}
\hfill
\begin{tabular}{rl}
\multicolumn{2}{l}{\textbf{PLUSQUAMPERFEKT}}\\
& hatten + \textbf{reduziert}\\
\end{tabular}

\vspace{10pt}

\begin{tabular}{lll}
\multicolumn{3}{l}{\textbf{MIT PRÄPOSITIONEN}}\\
& \textbf{reduzieren} + Akk & (to reduce something)\\
& \textbf{reduzieren auf} + Akk & (to reduce to)\\
& \textbf{reduzieren um} + Akk & (to reduce by)\\
\end{tabular}
\hfill
\begin{tabular}{lll}
\multicolumn{3}{l}{\textbf{TRENNBARE VERBEN}}\\
& \textbf{---} & \\
\end{tabular}

\vspace{-5pt}
\end{verbentry}

% registrieren (to register)
\begin{verbentry}{
  style = {default},
  toc = {registrieren (to register)},
  name = {registrieren},
  english = {to register},
  type = {regular, -ieren verb},
  auxiliary = {haben}
}
 
    
     \vspace{5pt}

    \hrule
    
    \vspace{5pt}

  \begin{tabular}{rl}
     \multicolumn{2}{l}{\textbf{PRÄSENS} }\\
    ich & \textbf{registriere}\\
    du & \textbf{registrierst}\\
    er/sie/es & \textbf{registriert}\\
    wir & \textbf{registrieren}\\
    ihr & \textbf{registriert}\\
    sie/Sie & \textbf{registrieren}\\
  \end{tabular}
  \hfill
     \begin{tabular}{rl}
    \multicolumn{2}{l}{\textbf{PRÄTERITUM} }\\
    ich & \textbf{registrierte}\\
    du & \textbf{registriertest}\\
    er/sie/es & \textbf{registrierte}\\
    wir & \textbf{registrierten}\\
    ihr & \textbf{registriertet}\\
    sie/Sie & \textbf{registrierten}\\
  \end{tabular}
  \hfill
  \begin{tabular}{rl}
     \multicolumn{2}{l}{\textbf{IMPERATIV} }\\
   & \textbf{registriere (du)!}\\
   & \textbf{registrieren wir!}\\
   & \textbf{registriert (ihr)!}\\
   & \textbf{registrieren Sie!}\\
   & \vspace{10pt}
  \end{tabular}

    \vspace{10pt}

 \begin{tabular}{rl}
     \multicolumn{2}{l}{\textbf{PERFEKT}}\\
   & haben + \textbf{registriert}\\
  \end{tabular}
\hfill
   \begin{tabular}{rl}
    \multicolumn{2}{l}{\textbf{FUTURE I} }\\
   & werden + \textbf{registrieren}\\
  \end{tabular}
\hfill
   \begin{tabular}{rl}
      \multicolumn{2}{l}{\textbf{PLUSQUAMPERFEKT}}\\
   & hatten + \textbf{registriert}\\
  \end{tabular}

   \vspace{10pt}

   \begin{tabular}{lll}
      \multicolumn{3}{l}{\textbf{MIT PRÄPOSITIONEN}}\\
& \textbf{sich registrieren bei} + Dat & (to register with)\\
& \textbf{sich registrieren für} + Akk & (to register for)\\
\end{tabular}
  \hfill
   \begin{tabular}{lll}
\multicolumn{3}{l}{\textbf{TRENNBARE VERBEN}}\\
& \textbf{---} & \\
\end{tabular}

  \vspace{-5pt}
\end{verbentry}

% regeln (to regulate / manage)
\begin{verbentry}{
  style = {default},
  toc = {regeln (to regulate / manage)},
  name = {regeln},
  english = {to regulate / manage},
  type = {regular},
  auxiliary = {haben}
}
 

\vspace{5pt}
\hrule
\vspace{5pt}

\begin{tabular}{rl}
\multicolumn{2}{l}{\textbf{PRÄSENS}}\\
ich & \textbf{regle}\\
du & \textbf{regelst}\\
er/sie/es & \textbf{regelt}\\
wir & \textbf{regeln}\\
ihr & \textbf{regelt}\\
sie/Sie & \textbf{regeln}\\
\end{tabular}
\hfill
\begin{tabular}{rl}
\multicolumn{2}{l}{\textbf{PRÄTERITUM}}\\
ich & \textbf{regelte}\\
du & \textbf{regeltest}\\
er/sie/es & \textbf{regelte}\\
wir & \textbf{regelten}\\
ihr & \textbf{regeltet}\\
sie/Sie & \textbf{regelten}\\
\end{tabular}
\hfill
\begin{tabular}{rl}
\multicolumn{2}{l}{\textbf{IMPERATIV}}\\
& \textbf{regle (du)!}\\
& \textbf{regeln wir!}\\
& \textbf{regelt (ihr)!}\\
& \textbf{regeln Sie!}\\
\end{tabular}

\vspace{10pt}

\begin{tabular}{rl}
\multicolumn{2}{l}{\textbf{PERFEKT}}\\
& haben + \textbf{geregelt}\\
\end{tabular}
\hfill
\begin{tabular}{rl}
\multicolumn{2}{l}{\textbf{FUTURE I}}\\
& werden + \textbf{regeln}\\
\end{tabular}
\hfill
\begin{tabular}{rl}
\multicolumn{2}{l}{\textbf{PLUSQUAMPERFEKT}}\\
& hatten + \textbf{geregelt}\\
\end{tabular}

\vspace{10pt}

\begin{tabular}{lll}
\multicolumn{3}{l}{\textbf{MIT PRÄPOSITIONEN}}\\
& \textbf{regeln} + Akk & (to manage something)\\
\end{tabular}
\hfill
\begin{tabular}{lll}
\multicolumn{3}{l}{\textbf{TRENNBARE VERBEN}}\\
& \textbf{---} & \\
\end{tabular}

\vspace{-5pt}
\end{verbentry}

% regnen (to rain)
\begin{verbentry}{
  style = {default},
  toc = {regnen (to rain)},
  name = {regnen},
  english = {to rain},
  type = {regular},
  auxiliary = {haben}
}
 
    
     \vspace{5pt}

    \hrule
    
    \vspace{5pt}

  \begin{tabular}{rl}
     \multicolumn{2}{l}{\textbf{PRÄSENS} }\\
    ich & \textbf{-}\\
    du & \textbf{-}\\
    es & \textbf{regnet}\\
    wir & \textbf{-}\\
    ihr & \textbf{-}\\
    sie/Sie & \textbf{-}\\
  \end{tabular}
  \hfill
     \begin{tabular}{rl}
    \multicolumn{2}{l}{\textbf{PRÄTERITUM} }\\
    ich & \textbf{-}\\
    du & \textbf{-}\\
    es & \textbf{regnete}\\
    wir & \textbf{-}\\
    ihr & \textbf{-}\\
    sie/Sie & \textbf{-}\\
  \end{tabular}
  \hfill
  \begin{tabular}{rl}
     \multicolumn{2}{l}{\textbf{IMPERATIV} }\\
   & \textbf{-}\\
   & \textbf{-}\\
   & \textbf{-}\\
   & \textbf{-}\\
   & \vspace{10pt}
  \end{tabular}

    \vspace{10pt}

 \begin{tabular}{rl}
     \multicolumn{2}{l}{\textbf{PERFEKT}}\\
  &  haben + \textbf{geredet}\\
  \end{tabular}
\hfill
   \begin{tabular}{rl}
   
    \multicolumn{2}{l}{\textbf{FUTURE I} }\\
  &  werden + \textbf{regnen}\\
 
  \end{tabular}
\hfill
   \begin{tabular}{rl}
      \multicolumn{2}{l}{\textbf{PLUSQUAMPERFEKT}}\\
  &  hatten + \textbf{geregnet}\\
  \end{tabular}
  
  \vspace{10pt}



\begin{tabular}{lll}
\multicolumn{3}{l}{\textbf{MIT PRÄPOSITIONEN}}\\
& \textbf{---} & \\
\end{tabular}
\hfill
\begin{tabular}{lll}
\multicolumn{3}{l}{\textbf{TRENNBARE VERBEN}}\\
& \textbf{---} & \\
\end{tabular}

  \vspace{-5pt}
\end{verbentry}

% reichen (to be enough / to hand)
\begin{verbentry}{
  style = {default},
  toc = {reichen (to be enough / to hand)},
  name = {reichen},
  english = {to be enough / to hand},
  type = {regular},
  auxiliary = {haben}
}
 

\vspace{5pt}
\hrule
\vspace{5pt}

\begin{tabular}{rl}
\multicolumn{2}{l}{\textbf{PRÄSENS}}\\
ich & \textbf{reiche}\\
du & \textbf{reichst}\\
er/sie/es & \textbf{reicht}\\
wir & \textbf{reichen}\\
ihr & \textbf{reicht}\\
sie/Sie & \textbf{reichen}\\
\end{tabular}
\hfill
\begin{tabular}{rl}
\multicolumn{2}{l}{\textbf{PRÄTERITUM}}\\
ich & \textbf{reichte}\\
du & \textbf{reichtest}\\
er/sie/es & \textbf{reichte}\\
wir & \textbf{reichten}\\
ihr & \textbf{reichtet}\\
sie/Sie & \textbf{reichten}\\
\end{tabular}
\hfill
\begin{tabular}{rl}
\multicolumn{2}{l}{\textbf{IMPERATIV}}\\
& \textbf{reiche (du)!}\\
& \textbf{reichen wir!}\\
& \textbf{reicht (ihr)!}\\
& \textbf{reichen Sie!}\\
\end{tabular}

\vspace{10pt}

\begin{tabular}{rl}
\multicolumn{2}{l}{\textbf{PERFEKT}}\\
& haben + \textbf{gereicht}\\
\end{tabular}
\hfill
\begin{tabular}{rl}
\multicolumn{2}{l}{\textbf{FUTURE I}}\\
& werden + \textbf{reichen}\\
\end{tabular}
\hfill
\begin{tabular}{rl}
\multicolumn{2}{l}{\textbf{PLUSQUAMPERFEKT}}\\
& hatten + \textbf{gereicht}\\
\end{tabular}

\vspace{10pt}

\begin{tabular}{lll}
\multicolumn{3}{l}{\textbf{MIT PRÄPOSITIONEN}}\\
& \textbf{reichen für} + Akk & (to be enough for)\\
& \textbf{reichen an} + Akk & (to hand to)\\
\end{tabular}
\hfill
\begin{tabular}{lll}
\multicolumn{3}{l}{\textbf{TRENNBARE VERBEN}}\\
& \textbf{ein$|$reichen} & (to submit)\\
& \textbf{nach$|$reichen} & (to hand later)\\
\end{tabular}

\vspace{-5pt}
\end{verbentry}

% reinigen (to clean)
\begin{verbentry}{
  style = {default},
  toc = {reinigen (to clean)},
  name = {reinigen},
  english = {to clean},
  type = {regular},
  auxiliary = {haben}
}


\vspace{5pt}
\hrule
\vspace{5pt}

\begin{tabular}{rl}
\multicolumn{2}{l}{\textbf{PRÄSENS}}\\
ich & \textbf{reinige}\\
du & \textbf{reinigst}\\
er/sie/es & \textbf{reinigt}\\
wir & \textbf{reinigen}\\
ihr & \textbf{reinigt}\\
sie/Sie & \textbf{reinigen}\\
\end{tabular}
\hfill
\begin{tabular}{rl}
\multicolumn{2}{l}{\textbf{PRÄTERITUM}}\\
ich & \textbf{reinigte}\\
du & \textbf{reinigtest}\\
er/sie/es & \textbf{reinigte}\\
wir & \textbf{reinigten}\\
ihr & \textbf{reinigtet}\\
sie/Sie & \textbf{reinigten}\\
\end{tabular}
\hfill
\begin{tabular}{rl}
\multicolumn{2}{l}{\textbf{IMPERATIV}}\\
& \textbf{reinige (du)!}\\
& \textbf{reinigen wir!}\\
& \textbf{reinigt (ihr)!}\\
& \textbf{reinigen Sie!}\\
\end{tabular}

\vspace{10pt}

\begin{tabular}{rl}
\multicolumn{2}{l}{\textbf{PERFEKT}}\\
& haben + \textbf{gereinigt}\\
\end{tabular}
\hfill
\begin{tabular}{rl}
\multicolumn{2}{l}{\textbf{FUTURE I}}\\
& werden + \textbf{reinigen}\\
\end{tabular}
\hfill
\begin{tabular}{rl}
\multicolumn{2}{l}{\textbf{PLUSQUAMPERFEKT}}\\
& hatten + \textbf{gereinigt}\\
\end{tabular}

\vspace{10pt}

\begin{tabular}{lll}
\multicolumn{3}{l}{\textbf{MIT PRÄPOSITIONEN}}\\
& \textbf{reinigen von} + Dat & (to clean from)\\
& \textbf{reinigen mit} + Dat & (to clean with)\\
\end{tabular}
\hfill
\begin{tabular}{lll}
\multicolumn{3}{l}{\textbf{TRENNBARE VERBEN}}\\
& \textbf{ab$|$reinigen} & (to clean off)\\
& \textbf{aus$|$reinigen} & (to clean out)\\
\end{tabular}
\vspace{-5pt}
\end{verbentry}

% reisen (to travel)
\begin{verbentry}{
  style = {default},
  toc = {reisen (to travel)},
  name = {reisen},
  english = {to travel},
  type = {regular},
  auxiliary = {sein}
}
 
    
     \vspace{5pt}

    \hrule
    
    \vspace{5pt}

  \begin{tabular}{rl}
     \multicolumn{2}{l}{\textbf{PRÄSENS} }\\
    ich & \textbf{reise}\\
    du & \textbf{reist}\\
    er/sie/es & \textbf{reist}\\
    wir & \textbf{reisen}\\
    ihr & \textbf{reist}\\
    sie/Sie & \textbf{reisen}\\
  \end{tabular}
  \hfill
     \begin{tabular}{rl}
    \multicolumn{2}{l}{\textbf{PRÄTERITUM} }\\
    ich & \textbf{reiste}\\
    du & \textbf{reistest}\\
    er/sie/es & \textbf{reiste}\\
    wir & \textbf{reisten}\\
    ihr & \textbf{reistet}\\
    sie/Sie & \textbf{reisten}\\
  \end{tabular}
  \hfill
  \begin{tabular}{rl}
     \multicolumn{2}{l}{\textbf{IMPERATIV} }\\
   & \textbf{reis(e) (du)!}\\
   & \textbf{reisen wir!}\\
   & \textbf{reist (ihr)!}\\
   & \textbf{reisen Sie!}\\
   & \vspace{10pt}
  \end{tabular}

    \vspace{10pt}

 \begin{tabular}{rl}
     \multicolumn{2}{l}{\textbf{PERFEKT}}\\
   & sein + \textbf{gereist}\\
  \end{tabular}
\hfill
   \begin{tabular}{rl}
   
    \multicolumn{2}{l}{\textbf{FUTURE I} }\\
  &  werden + \textbf{reisen}\\
 
  \end{tabular}
\hfill
   \begin{tabular}{rl}
      \multicolumn{2}{l}{\textbf{PLUSQUAMPERFEKT}}\\
  &  waren + \textbf{gereist}\\
  \end{tabular}
  
  \vspace{10pt}



\begin{tabular}{lll}
\multicolumn{3}{l}{\textbf{MIT PRÄPOSITIONEN}}\\
& \textbf{---} & \\
\end{tabular}
\hfill
\begin{tabular}{lll}
\multicolumn{3}{l}{\textbf{TRENNBARE VERBEN}}\\
& \textbf{---} & \\
\end{tabular}

  \vspace{-5pt}
\end{verbentry}

% reklamieren (to complain / claim)
\begin{verbentry}{
  style = {default},
  toc = {reklamieren (to complain / claim)},
  name = {reklamieren},
  english = {to complain / claim},
  type = {regular},
  auxiliary = {haben}
}


\vspace{5pt}
\hrule
\vspace{5pt}

\begin{tabular}{rl}
\multicolumn{2}{l}{\textbf{PRÄSENS}}\\
ich & \textbf{reklamiere}\\
du & \textbf{reklamierst}\\
er/sie/es & \textbf{reklamiert}\\
wir & \textbf{reklamieren}\\
ihr & \textbf{reklamiert}\\
sie/Sie & \textbf{reklamieren}\\
\end{tabular}
\hfill
\begin{tabular}{rl}
\multicolumn{2}{l}{\textbf{PRÄTERITUM}}\\
ich & \textbf{reklamierte}\\
du & \textbf{reklamiertest}\\
er/sie/es & \textbf{reklamierte}\\
wir & \textbf{reklamierten}\\
ihr & \textbf{reklamiertet}\\
sie/Sie & \textbf{reklamierten}\\
\end{tabular}
\hfill
\begin{tabular}{rl}
\multicolumn{2}{l}{\textbf{IMPERATIV}}\\
& \textbf{reklamiere (du)!}\\
& \textbf{reklamieren wir!}\\
& \textbf{reklamiert (ihr)!}\\
& \textbf{reklamieren Sie!}\\
\end{tabular}

\vspace{10pt}

\begin{tabular}{rl}
\multicolumn{2}{l}{\textbf{PERFEKT}}\\
& haben + \textbf{reklamiert}\\
\end{tabular}
\hfill
\begin{tabular}{rl}
\multicolumn{2}{l}{\textbf{FUTURE I}}\\
& werden + \textbf{reklamieren}\\
\end{tabular}
\hfill
\begin{tabular}{rl}
\multicolumn{2}{l}{\textbf{PLUSQUAMPERFEKT}}\\
& hatten + \textbf{reklamiert}\\
\end{tabular}

\vspace{10pt}

\begin{tabular}{lll}
\multicolumn{3}{l}{\textbf{MIT PRÄPOSITIONEN}}\\
& \textbf{reklamieren bei} + Dat & (to complain to)\\
& \textbf{reklamieren für} + Akk & (to claim for)\\
\end{tabular}
\hfill
\begin{tabular}{lll}
\multicolumn{3}{l}{\textbf{TRENNBARE VERBEN}}\\
& \textbf{---} & \\
\end{tabular}
\vspace{-5pt}
\end{verbentry}

% rufen (to call / shout)
\begin{verbentry}{
  style = {acc},
  toc = {rufen (to call / shout)},
  name = {rufen},
  english = {to call / shout},
  type = {irregular, + Akkusativ},
  auxiliary = {haben}
}
 
    
     \vspace{5pt}

    \hrule
    
    \vspace{5pt}

  \begin{tabular}{rl}
     \multicolumn{2}{l}{\textbf{PRÄSENS} }\\
    ich & \textbf{rufe}\\
    du & \textbf{rufst}\\
    er/sie/es & \textbf{ruft}\\
    wir & \textbf{rufen}\\
    ihr & \textbf{ruft}\\
    sie/Sie & \textbf{rufen}\\
  \end{tabular}
  \hfill
     \begin{tabular}{rl}
    \multicolumn{2}{l}{\textbf{PRÄTERITUM} }\\
    ich & \textbf{rief}\\
    du & \textbf{riefst}\\
    er/sie/es & \textbf{rief}\\
    wir & \textbf{riefen}\\
    ihr & \textbf{rieft}\\
    sie/Sie & \textbf{riefen}\\
  \end{tabular}
  \hfill
  \begin{tabular}{rl}
     \multicolumn{2}{l}{\textbf{IMPERATIV} }\\
   & \textbf{ruf(e) (du)!}\\
   & \textbf{rufen wir!}\\
   & \textbf{ruft (ihr)!}\\
   & \textbf{rufen Sie!}\\
   & \vspace{10pt}
  \end{tabular}

    \vspace{10pt}

 \begin{tabular}{rl}
     \multicolumn{2}{l}{\textbf{PERFEKT}}\\
   & haben + \textbf{gerufen}\\
  \end{tabular}
\hfill
   \begin{tabular}{rl}
   
    \multicolumn{2}{l}{\textbf{FUTURE I} }\\
   & werden + \textbf{rufen}\\
 
  \end{tabular}
\hfill
   \begin{tabular}{rl}
      \multicolumn{2}{l}{\textbf{PLUSQUAMPERFEKT}}\\
   & hatten + \textbf{gerufen}\\
  \end{tabular}
  

 \vspace{10pt}

   \begin{tabular}{lll}
      \multicolumn{3}{l}{\textbf{MIT PRÄPOSITIONEN}}\\
& \textbf{rufen nach} + Dat & (to cry out for)\\
& \textbf{rufen zu} + Dat & (to call to)\\
&\textbf{rufen um} + Akk & (to plead for)\\
&\textbf{anrufen bei} + Dat & (to call at (office))\\
&\textbf{anrufen wegen} + Gen/Dat & (to call about)\\
&\textbf{zurufen zu} + Dat & (to shout to)\\
&\textbf{aufrufen zu} + Dat & (to call on to do)\\
&\textbf{abrufen von} + Dat & (to retrieve from)\\
  \end{tabular}
  \hfill
   \begin{tabular}{lll}
      \multicolumn{3}{l}{\textbf{TRENNBARE VERBEN}}\\
 & \textbf{an$|$rufen} & (to call (phone))\\
 & \textbf{ab$|$rufen} & (to retrieve / access)\\
 & \textbf{zu$|$rufen} & (to call out)\\
  & \textbf{heraus$|$rufen} & (to call out (from inside))\\
  & \textbf{hinein$|$rufen} & (to call into)\\
  & \textbf{zurück$|$rufen} & (to call back)\\
  & \textbf{auf$|$rufen} & (to call up / apppeal)\\
     & \textbf{herbei$|$rufen} & (to summon)\\  
  \end{tabular}


  \vspace{-5pt}
\end{verbentry}

% ruhen (to rest)
\begin{verbentry}{
  style = {default},
  toc = {ruhen (to rest)},
  name = {ruhen},
  english = {to rest},
  type = {regular},
  auxiliary = {haben}
}
 

\vspace{5pt}
\hrule
\vspace{5pt}

\begin{tabular}{rl}
\multicolumn{2}{l}{\textbf{PRÄSENS}}\\
ich & \textbf{ruhe}\\
du & \textbf{ruhst}\\
er/sie/es & \textbf{ruht}\\
wir & \textbf{ruhen}\\
ihr & \textbf{ruht}\\
sie/Sie & \textbf{ruhen}\\
\end{tabular}
\hfill
\begin{tabular}{rl}
\multicolumn{2}{l}{\textbf{PRÄTERITUM}}\\
ich & \textbf{ruhte}\\
du & \textbf{ruhtest}\\
er/sie/es & \textbf{ruhte}\\
wir & \textbf{ruhten}\\
ihr & \textbf{ruhtet}\\
sie/Sie & \textbf{ruhten}\\
\end{tabular}
\hfill
\begin{tabular}{rl}
\multicolumn{2}{l}{\textbf{IMPERATIV}}\\
& \textbf{ruhe (du)!}\\
& \textbf{ruhen wir!}\\
& \textbf{ruht (ihr)!}\\
& \textbf{ruhen Sie!}\\
\end{tabular}

\vspace{10pt}

\begin{tabular}{rl}
\multicolumn{2}{l}{\textbf{PERFEKT}}\\
& haben + \textbf{geruht}\\
\end{tabular}
\hfill
\begin{tabular}{rl}
\multicolumn{2}{l}{\textbf{FUTURE I}}\\
& werden + \textbf{ruhen}\\
\end{tabular}
\hfill
\begin{tabular}{rl}
\multicolumn{2}{l}{\textbf{PLUSQUAMPERFEKT}}\\
& hatten + \textbf{geruht}\\
\end{tabular}

\vspace{10pt}

\begin{tabular}{lll}
\multicolumn{3}{l}{\textbf{MIT PRÄPOSITIONEN}}\\
& \textbf{ruhen auf} + Dat & (to rest on)\\
\end{tabular}
\hfill
\begin{tabular}{lll}
\multicolumn{3}{l}{\textbf{TRENNBARE VERBEN}}\\
& \textbf{---} & \\
\end{tabular}

\vspace{-5pt}
\end{verbentry}
